\documentclass{article}

\usepackage{xparse}

\usepackage{xltxtra}

\usepackage{setspace}

\usepackage{fontspec}

\setmainfont[Path="V:/Fonts/Eng/",AutoFakeBold=3,AutoFakeSlant=0.25]{BASKVILL.TTF}
\setsansfont[Path="V:/Fonts/Eng/",AutoFakeBold=3,AutoFakeSlant=0.25]{POORICH.TTF}
\setmonofont[Path="V:/Fonts/Eng/",AutoFakeBold=3,AutoFakeSlant=0.25]{cour.ttf}

\usepackage{xeCJK}
\setCJKmainfont[Path="V:/Fonts/Zh/",AutoFakeSlant=0.2,AutoFakeBold=3]{STZHONGS.TTF}
\setCJKsansfont[Path="V:/Fonts/Zh/",AutoFakeSlant=0.2,AutoFakeBold=3]{STKAITI.ttf}
\setCJKmonofont[Path="V:/Fonts/Zh/",AutoFakeSlant=0.2,AutoFakeBold=3]{simfang.ttf}

\newfontfamily{\yh}[AutoFakeSlant=0.225,AutoFakeBold=3]{Microsoft YaHei}
\newfontfamily{\fontik}[AutoFakeSlant=0.225,AutoFakeBold=3,Path="V:/Fonts/Eng/"]{Inkfree.ttf}
\newfontfamily{\fontss}[AutoFakeSlant=0.225,AutoFakeBold=3,Path="V:/Fonts/Eng/"]{HARNGTON.TTF}
\newfontfamily{\fontfl}[AutoFakeSlant=0.225,AutoFakeBold=3,Path="V:/Fonts/Eng/"]{FTLTLT.TTF}
\newfontfamily{\fontch}[AutoFakeSlant=0.225,AutoFakeBold=3,Path="V:/Fonts/Eng/"]{CHILLER.TTF}
\newfontfamily{\fontrm}[AutoFakeSlant=0.225,AutoFakeBold=3,Path="V:/Fonts/Eng/"]{ERASLGHT.TTF}
\newfontfamily{\fontbk}[AutoFakeSlant=0.225,AutoFakeBold=3,Path="V:/Fonts/Eng/"]{BOOKOS.TTF}
\newfontfamily{\fontlr}[AutoFakeSlant=0.225,AutoFakeBold=3,Path="V:/Fonts/Eng/"]{LBRITEI.TTF}

\usepackage{amsmath,amssymb,amsthm,amsfonts,mathrsfs}
\usepackage{mathtools}
\usepackage{physics}
\usepackage{esint}
\usepackage{bm}

\usepackage{xcolor}

\usepackage{colortbl} % 单元格颜色
\usepackage{makecell} % 单元格换行
\usepackage{multirow} % 合并列
\usepackage{longtable} % 跨页长表格
\usepackage{booktabs} % 三线表

\definecolor{backcolor}{rgb}{0.95,0.95,0.92}
\definecolor{codegray}{gray}{0.6}

\definecolor{darkred}{RGB}{139,0,0}
\definecolor{darkpurple}{RGB}{143, 66, 197}

\definecolor{line15}{RGB}{200,179,142}

\definecolor{green1}{RGB}{0, 156, 39} % 深绿
\definecolor{green2}{RGB}{214, 254, 224} % 浅绿
\definecolor{blue1}{RGB}{163, 172, 204} % 深蓝
\definecolor{blue2}{RGB}{204, 252, 254} % 浅蓝

\definecolor{blue3}{HTML}{D1FCFE} % Theorem blue
\definecolor{blue4}{RGB}{72, 141, 246} % Definition blue
\definecolor{blue5}{HTML}{F0FEFF} % Lemma blue
\definecolor{blue6}{RGB}{20, 230, 253} % Criterion blue
\definecolor{blue7}{RGB}{207, 225, 235} % blue7

\usepackage[margin=1.5cm]{geometry}

\setlength{\parindent}{0pt}

\linespread{1.3}

\usepackage{minted}
\usemintedstyle{default}
\setminted{
    %linenos,
    breaklines,
}

\usepackage[most]{tcolorbox}

\newtcbox{\highl}[1][red]{on line,
    arc=0pt,outer arc=0pt,colback=#1!10!white,colframe=#1!50!black,
    boxsep=0pt,left=1pt,right=1pt,top=2pt,bottom=2pt,
    boxrule=0pt,bottomrule=1pt,
}

\newtcbox{\chighl}[1][red]{on line,
    arc=7pt,colback=#1!10!white,colframe=#1!50!black,
    before upper={\rule[-3pt]{0pt}{10pt}},boxrule=1pt,
    boxsep=0pt,left=6pt,right=6pt,top=2pt,bottom=2pt
}

% ========================= Envs ================================

% counters:
\newcounter{thmcounter}[subsection]
\setcounter{thmcounter}{0}
\newcounter{dfncounter}[subsection]
\setcounter{dfncounter}{0}
\newcounter{xmpcounter}[subsection]
\setcounter{xmpcounter}{0}
\newcounter{pcdcounter}
\setcounter{pcdcounter}{0}
\newcounter{notecounter}
\setcounter{notecounter}{0}
\newcommand{\thethmcountercustom}{\thesection.\arabic{thmcounter}}
\newcommand{\thedfncountercustom}{\thesection.\arabic{dfncounter}}
\newcommand{\thexmpcountercustom}{\thesection.\arabic{xmpcounter}}
\newcommand{\thethmcountercustomsub}{\thesubsection.\arabic{thmcounter}}
\newcommand{\thedfncountercustomsub}{\thesubsection.\arabic{dfncounter}}
\newcommand{\thexmpcountercustomsub}{\thesubsection.\arabic{xmpcounter}}

\NewDocumentEnvironment{qst}{s O{NaN}}{\begin{tcolorbox}
    [enhanced, colback=white, boxrule=0pt, frame hidden, borderline west={0.7mm}{0.1mm}{darkred}, breakable]
    \textbf{{\Large Question~#2}}\phantomsection\label{qst:#2}
    \IfBooleanTF{#1}{\newline}{}
}{\end{tcolorbox}}

% proof:
\newenvironment{prf}{\begin{tcolorbox}
    [enhanced, colback=lightgray!50!white, boxrule=0pt, frame hidden, breakable, after upper={\hfill$\square$}]
    {{\large\textit{Proof:}}\newline\fontik}
}{\end{tcolorbox}}

% remark:
\newenvironment{rmk}[1][contributor]{\begin{tcolorbox}
    [enhanced, colback=line15!20!white, boxrule=0pt, frame hidden, fontupper=\sffamily,
        borderline west={0.7mm}{0.1mm}{line15}, breakable, after upper={\hfill\textcolor{codegray}{\textit{remarked~by~#1}}}]
    {\textbf{{\large\rmfamily Remark}}}\newline
}{\end{tcolorbox}}

\newenvironment{intrormk}[1][contributor]{\begin{tcolorbox}
    [enhanced, colback=line15!10!white, boxrule=0pt, frame hidden, fontupper=\sffamily,
        borderline north={0.7mm}{0.1mm}{line15}, breakable, after upper={\hfill\textcolor{codegray}{\textit{remarked~by~#1}}}]
    {\textbf{{\centering\large\rmfamily Introduction}}}\newline
}{\end{tcolorbox}}

% definition:
\NewDocumentEnvironment{dfn}{s o m s o O{contributor}}{\begin{tcolorbox}
    [enhanced, colback=white, boxrule=0pt, frame hidden, borderline west={0.7mm}{0.1mm}{blue4}, breakable, after upper={\IfBooleanTF{#4}{\hfill\textcolor{codegray}{\textit{by~#6}}}{}}]
    {\stepcounter{dfncounter}\textbf{{\large Definition \thedfncountercustomsub:~#3}\IfNoValueTF{#5}{}{\IfBlankTF{#5}{}{(#5)}}}\IfNoValueTF{#2}{}{\phantomsection\label{dfn:#2}}}
    \IfBooleanTF{#1}{}{\newline}
}{\end{tcolorbox}}

\NewDocumentEnvironment{vardfn}{s o m s o O{contributor}}{\begin{tcolorbox}
    [enhanced, colback=white, boxrule=0pt, frame hidden, borderline west={0.7mm}{0.1mm}{blue4, dotted}, breakable, after upper={\IfBooleanTF{#4}{\hfill\textcolor{codegray}{\textit{by~#6}}}{}}]
    {\stepcounter{dfncounter}\textbf{{\large Definition \thedfncountercustomsub:~#3 *}\IfNoValueTF{#5}{}{\IfBlankTF{#5}{}{(#5)}}}\IfNoValueTF{#2}{}{\phantomsection\label{dfn:#2}}}
    \IfBooleanTF{#1}{}{\newline}
}{\end{tcolorbox}}

% theorem:
\NewDocumentEnvironment{thm}{s o m s o O{contributor}}{\begin{tcolorbox} % bool(newline), label, name, bool(contributor displayed), subname, contributor
    [enhanced, boxrule=0pt, frame hidden, borderline west={0.7mm}{0.1mm}{blue1}, colback=blue3, breakable, after upper={\IfBooleanTF{#4}{\hfill\textcolor{codegray}{\textit{by~#6}}}{}}]
    {\stepcounter{thmcounter}\textbf{{\large Theorem \thethmcountercustomsub:~#3}\IfNoValueTF{#5}{}{\IfBlankTF{#5}{}{(#5)}}}\IfNoValueTF{#2}{}{\phantomsection\label{thm:#2}}}
    \IfBooleanTF{#1}{}{\newline}
}{\end{tcolorbox}}

\NewDocumentEnvironment{varthm}{s o m s o O{contributor}}{\begin{tcolorbox} % bool(newline), label, name, bool(contributor displayed), subname, contributor
    [enhanced, boxrule=0pt, frame hidden, borderline west={0.7mm}{0.1mm}{blue1}, interior style={left color=blue1, right color=blue3, middle color=white}, sharp corners, breakable, after upper={\IfBooleanTF{#4}{\hfill\textcolor{codegray}{\textit{by~#6}}}{}}]
    {\stepcounter{thmcounter}\textbf{{\large Theorem \thethmcountercustomsub:~#3 *}\IfNoValueTF{#5}{}{\IfBlankTF{#5}{}{(#5)}}}\IfNoValueTF{#2}{}{\phantomsection\label{thm:#2}}}
    \IfBooleanTF{#1}{}{\newline}
}{\end{tcolorbox}}

% property theorem and criterion theorem 
\NewDocumentEnvironment{ppt}{s o m s o O{contributor}}{\begin{tcolorbox} % star, label, name, subname, contributor
    [enhanced, colback=white, boxrule=0pt, frame hidden, borderline west={0.7mm}{0.1mm}{blue1}, breakable, after upper={\IfBooleanTF{#4}{\hfill\textcolor{codegray}{\textit{by~#6}}}{}}]
    {\stepcounter{thmcounter}\textbf{{Property Theorem} \thethmcountercustomsub~#3\IfNoValueTF{#5}{}{\IfBlankTF{#5}{}{(#5)}}:}\IfNoValueTF{#2}{}{\phantomsection\label{ppt:#2}}}
    \IfBooleanTF{#1}{\newline}{}
}{\end{tcolorbox}}

\NewDocumentEnvironment{crt}{s o m s o O{contributor}}{\begin{tcolorbox} % star, label, name, subname, contributor
    [enhanced, colback=white, boxrule=0pt, frame hidden, 
    borderline west={0.7mm}{0.1mm}{blue6}, interior style={left color=white, right color=blue6!50!white}, 
    breakable, after upper={\IfBooleanTF{#4}{\hfill\textcolor{codegray}{\textit{by~#6}}}{}}]
    {\stepcounter{thmcounter}\textbf{{Criterion Theorem} \thethmcountercustomsub~#3\IfNoValueTF{#5}{}{\IfBlankTF{#5}{}{(#5)}}:}\IfNoValueTF{#2}{}{\phantomsection\label{thm:#2}}}
    \IfBooleanTF{#1}{\newline}{}
}{\end{tcolorbox}}

% corollary:
\NewDocumentEnvironment{crl}{s o m s o O{contributor}}{\begin{tcolorbox} % star, label, name, subname, contributor
    [enhanced, colback=white!95!blue, boxrule=1pt, frame hidden, 
    borderline west={0.7mm}{0.1mm}{blue1}, 
    breakable, after upper={\IfBooleanTF{#4}{\hfill\textcolor{codegray}{\textit{by~#6}}}{}}]
    {\stepcounter{thmcounter}\textbf{{Corollary} \thethmcountercustomsub~#3\IfNoValueTF{#5}{}{\IfBlankTF{#5}{}{(#5)}}:}\IfNoValueTF{#2}{}{\phantomsection\label{thm:#2}}}
    \IfBooleanTF{#1}{}{\newline}
}{\end{tcolorbox}}

% exmaple and counterexample:
\NewDocumentEnvironment{xmp}{s o m s o O{contributor}}{\begin{tcolorbox} % star, label, name, subname, contributor
    [enhanced, colback=white, boxrule=0pt, frame hidden, borderline west={0.7mm}{0.1mm}{green1}, breakable, after upper={\IfBooleanTF{#4}{\hfill\textcolor{codegray}{\textit{by~#6}}}{}}]
    {\stepcounter{xmpcounter}\textbf{{Example} \thexmpcountercustomsub~#3\IfNoValueTF{#5}{}{\IfBlankTF{#5}{}{(#5)}}:}\IfNoValueTF{#2}{}{\phantomsection\label{xmp:#2}}}
    \IfBooleanTF{#1}{}{\newline}
}{\end{tcolorbox}}

\NewDocumentEnvironment{cxmp}{s o m s o O{contributor}}{\begin{tcolorbox} % star, label, name, subname, contributor
    [enhanced, colback=green1!25!white, boxrule=0pt, frame hidden, 
    borderline west={0.7mm}{0.1mm}{green1}, breakable, after upper={\IfBooleanTF{#4}{\hfill\textcolor{codegray}{\textit{by~#6}}}{}}]
    {\stepcounter{xmpcounter}\textbf{{Counterexample} \thexmpcountercustomsub~#3\IfNoValueTF{#5}{}{\IfBlankTF{#5}{}{(#5)}}:}\IfNoValueTF{#2}{}{\phantomsection\label{cxmp:#2}}}
    \IfBooleanTF{#1}{}{\newline}
}{\end{tcolorbox}}

% instance:
\NewDocumentEnvironment{ins}{s o m s o O{contributor}}{\begin{tcolorbox}
    [enhanced, colback=green2!50!white, boxrule=0pt, frame hidden, borderline west={0.7mm}{0.1mm}{green1}, breakable, after upper={\IfBooleanTF{#4}{\hfill\textcolor{codegray}{\textit{by~#6}}}{}}]
    {\stepcounter{xmpcounter}\textbf{{Instance} \thexmpcountercustomsub~~#3\IfNoValueTF{#5}{}{\IfBlankTF{#5}{}{(#5)}}:}\IfNoValueTF{#2}{}{\phantomsection\label{xmp:#2}}}
    \IfBooleanTF{#1}{\newline}{}
}{\end{tcolorbox}}

% ========================Reused Envs==========================

\NewDocumentEnvironment{note}{o m o}{\begin{tcolorbox}
    [enhanced, coltext=black, colback=white!95!yellow, colframe=black, boxrule=0.4mm, arc=1mm, 
    before upper={\hspace{2mm}\faQuoteLeft~}, after upper={~\faQuoteRight}, breakable, 
    title=Note \thenotecounter : #2\IfNoValueTF{#1}{}{\phantomsection\label{note:#1}}, fonttitle=\bfseries\large, left=1em, right=1em, top=0.5em, bottom=0.5em]
    {\stepcounter{notecounter}\IfNoValueTF{#3}{}{\IfBlankTF{#3}{}{\textbf{#3}:\newline}}}
}{\end{tcolorbox}}

\NewDocumentEnvironment{varnote}{o m o}{\begin{tcolorbox}
    [enhanced, coltext=black, colback=white!95!blue, colframe=blue1, boxrule=0.4mm, arc=1mm, 
    before upper={\hspace{2mm}\faQuoteLeft~}, after upper={~\faQuoteRight}, breakable, 
    title=Note \thenotecounter : #2\IfNoValueTF{#1}{}{\phantomsection\label{note:#1}}, fonttitle=\bfseries\large, left=1em, right=1em, top=0.5em, bottom=0.5em]
    {\stepcounter{notecounter}\IfNoValueTF{#3}{}{\IfBlankTF{#3}{}{\textbf{#3}\newline}}}
}{\end{tcolorbox}}

\usepackage{graphicx}

\usepackage{subcaption}

\usepackage{enumitem}

\usepackage{hyperref}

\hypersetup{
    colorlinks=true,
    linkcolor=blue,
    filecolor=magenta,
    urlcolor=cyan
}

% Tikz
\usepackage{tikz}
\usetikzlibrary{shapes,arrows,positioning,calc,decorations.pathreplacing,patterns,3d,intersections}

% 自定义的数学运算符
\DeclareMathOperator{\N}{\mathbb{N}}                    % Set of Natural Numbers
\DeclareMathOperator{\Nat}{\mathsf{Nat}}
\DeclareMathOperator{\Z}{\mathbb{Z}}                    % Set of Integers
\DeclareMathOperator{\Q}{\mathbb{Q}}                    % Set of Rational Numbers
\DeclareMathOperator{\R}{\mathbb{R}}                    % Set of Real Numbers
\DeclareMathOperator{\C}{\mathbb{C}}                    % Set of Complex Numbers
\DeclareMathOperator{\D}{\mathbb{D}}
\DeclareMathOperator{\0}{\mathbf{0}}
\DeclareMathOperator{\x}{\mathbf{x}}
\DeclareMathOperator{\w}{\mathbf{w}}
\DeclareMathOperator{\ii}{\mathrm{i}}
\DeclareMathOperator{\st}{\text{s.t.}}
\DeclareMathOperator{\B}{\mathbf{B}}
\DeclareMathOperator{\Ln}{Ln}
\DeclareMathOperator{\Arg}{Arg}
\DeclareMathOperator{\Log}{Log}

\DeclareMathOperator{\card}{card}
\DeclareMathOperator{\dist}{dist}

\DeclareMathOperator{\Id}{Id}                           % Identity
\DeclareMathOperator{\img}{im}                          % Image

\DeclareMathOperator{\Hom}{Hom}                         % Set of Homomorphisms
\DeclareMathOperator{\End}{End}                         % Set of Endomorphisms
\DeclareMathOperator{\Aut}{Aut}                         % Set of Automorphisms
\DeclareMathOperator{\Isom}{Isom}                       % Set of Isomorphisms

\DeclareMathOperator{\diag}{diag}                       % Diagonal Matrix

\DeclareMathOperator{\F}{\mathbb{F}}                    % Number Field (F)
\DeclareMathOperator{\mlim}{\xrightarrow{m}}
\DeclareMathOperator{\aelim}{\xrightarrow{\text{a.e.}}}

% Operators
\DeclareMathOperator{\blo}{\mathscr{B}}

% Hyperrefs

% Colored:
\newcommand{\colora}[1]{\textcolor{red}{#1}}
\newcommand{\colorb}[1]{\textcolor{blue4}{#1}}
\newcommand{\colorc}[1]{\textcolor{blue6}{#1}}
\newcommand{\colord}[1]{\textcolor{darkpurple}{#1}}

% Logic:
\newcommand{\tru}{\textcolor{green1}{\texttt{True}}}
\newcommand{\fls}{\textcolor{red}{\texttt{False}}}

% Non-mathematical
\newcommand{\nohref}[1]{\textcolor{darkpurple}{``#1''}}
\newcommand{\nohrefconcept}[1]{\textcolor{blue4}{``#1''}}

\newcommand{\continued}{{\Large\textbf{To be Continued...}}}
\newcommand{\bysorry}{\fcolorbox{white}{gray!20!white}{\texttt{\textcolor{black}{:= }\textcolor{blue}{by }\textcolor{red}{\textbf{sorry}}}}\;\faTimesCircle\;\,\texttt{\textcolor{gray}{{\footnotesize declaration uses 'sorry'.}}}}

% signature
\newcommand{\Dedicatia}{{\fontik\textcolor{darkpurple}{Dedicatia}}}

\usepackage{fancyhdr}
\pagestyle{fancy}

\fancyhf{}

\renewcommand{\headrulewidth}{0.4pt}
\renewcommand{\footrulewidth}{0.25pt}

\fancyhead[L]{Some Notes about Functional Analysis}
\fancyhead[R]{MATH3611-01}
\fancyfoot[C]{Page.\thepage}

\title{Some Notes about Functional Analysis}
\author{Dedicatiaaaaa}
\date{Mar. 3rd, 2026}

% Preamble ends here

\begin{document}
\maketitle
\vspace{4em}
\begin{center}
    This document is made by \XeLaTeX{}.\\
    \vspace{\baselineskip}
    %{\textcolor{red}{Due: Dec. 30th, 2025}}\hspace{2em} 
    {\textcolor{darkpurple}{Last Commit: \today}} 
\end{center}
\newpage
\tableofcontents
\newpage
\section{基本定义}

    \subsection{Banach 空间}

        \begin{dfn}
            [Norm]
            {范数}
            [Norm]
            数域 \(\F\) 上的线性空间 $X$ 上的一个函数 $\norm{\cdot}:X\to\R$ 是 $X$ 上的一个\textbf{范数}, 当且仅当:
            \begin{enumerate}[label=(\roman*)]
                \item \textbf{非退化性 / 正定性}: $\norm{x}=0\iff x=0$;
                \item \textbf{绝对齐次性}: $\norm{\alpha x}=\abs{\alpha}\norm{x}$, $\forall \alpha\in\F, x\in X$;
                \item \textbf{次可加性 / 三角不等式}: $\norm{x+y}\leq\norm{x}+\norm{y}$, $\forall x,y\in X$.
            \end{enumerate}
        \end{dfn}

        \begin{dfn}
            [NVS]
            {赋范线性空间}
            [Normed Vector Space]
            线性空间 \(X\) 是赋范线性空间当且仅当 \(X\) 上存在范数函数.
        \end{dfn}

        \begin{rmk}
            [Dedicatia]
            一般称 \((X, \norm{\cdot})\) 是一个赋范线性空间, 其中 \(X\) 是线性空间, \(\norm{\cdot}\) 是 \(X\) 上的范数函数. 也称 \(X\) 是配备了范数 \(\norm{\cdot}\) 的赋范线性空间.
        \end{rmk}

        \begin{ppt}
            {赋范线性空间是度量空间}
            [Normed Vector Space is a Metric Space]
            范数可以诱导度量: 具体来说, 对于赋范线性空间 \((X, \norm{\cdot})\), \(x,y\in X\), 定义
            \[d(x,y)=\norm{x-y}\]
            是 \(X\) 上的一个度量函数, 从而 \(X\) 也是一个度量空间.
        \end{ppt}

        \begin{rmk}
            [Dedicatia]
            一般的度量不一定能诱导范数.
        \end{rmk}

        \begin{dfn}
            [BS]
            {Banach 空间}
            [Banach Space]
            赋范线性空间 \((X, \norm{\cdot})\) 是 Banach 空间当且仅当 \(X\) 关于由范数诱导的度量 \(d(x,y)=\norm{x-y}\) 是完备的.
        \end{dfn}

        \begin{rmk}
            [Dedicatia]
            所谓\textbf{完备}, 就是 \(X\) 中的每个 Cauchy 列都收敛于 \(X\) 中的某个元素.
        \end{rmk}

        \continued

    \subsection{线性算子}
    
        \begin{dfn}
            [BLO]
            {有界线性算子}
            [Bounded Linear Operator]
            设 \((X, \norm{\cdot}_X)\) 和 \((Y, \norm{\cdot}_Y)\) 是两个赋范线性空间. 定义线性映射 \(T:X\to Y\) 是一个\textbf{有界线性算子}, 当且仅当存在 \(M\geq 0\) 使得
            \[\forall x\in X, \norm{Tx}_Y\leq M\norm{x}_X\]
        \end{dfn}

        \begin{dfn}
            [BLF]
            {有界线性泛函}
            [Bounded Linear Functional]
            设 \((X, \norm{\cdot})\) 是一个赋范线性空间. 定义线性映射 \(f:X\to\R\) 是一个\textbf{有界线性泛函}, 当且 \(f\) 是一个有界线性算子.
        \end{dfn}

        \begin{rmk}
            [Dedicatia]
            直观地理解, 有界线性算子将 \(X\) 的单位闭球映到 \(Y\) 中半径不超过 \(M\) 的单位闭球中.
        \end{rmk}

        \begin{dfn}
            {有界线性算子的范数}
            [Norm of Bounded Linear Operator]
            设 \((X, \norm{\cdot}_X)\) 和 \((Y, \norm{\cdot}_Y)\) 是两个赋范线性空间, \(T:X\to Y\) 是一个有界线性算子. 定义 \(T\) 的范数为
            \[\norm{T}:=\inf\{M\geq 0 : \forall x\in X, \norm{Tx}_Y\leq M\norm{x}_X\}\]
        \end{dfn}

        \begin{rmk}
            [Dedicatia]
            \(T\) 的范数 \(\norm{T}\) 是满足 \(\forall x\in X, \norm{Tx}_Y\leq M\norm{x}_X\) 的最小的 \(M\).
        \end{rmk}

        \begin{ppt}
            {有界线性算子范数的等价定义}
            [Equivalent Definition of Norm of Bounded Linear Operator]
            设 \((X, \norm{\cdot}_X)\) 和 \((Y, \norm{\cdot}_Y)\) 是两个赋范线性空间, \(T:X\to Y\) 是一个有界线性算子. 则 \(T\) 的范数为
            \[\norm{T}=\sup_{x\neq 0}\frac{\norm{Tx}_Y}{\norm{x}_X}=\sup_{\norm{x}_X=1}\norm{Tx}_Y\]
        \end{ppt}

        \begin{rmk}
            [Dedicatia]
            为什么可以这样等价, 是因为 \(X,Y\) 都是线性空间, 可以进行标量乘法. \\
            所以对于任意 \(x\in X\), 当 \(x\neq 0\) 时, 可以将 \(x\) 写成 \(\norm{x}_X\cdot\frac{x}{\norm{x}_X}\). 从而 \(\norm{Tx}_Y=\norm{T(\norm{x}_X\cdot\frac{x}{\norm{x}_X})}_Y=\norm{x}_X\norm{T(\frac{x}{\norm{x}_X})}_Y\). 因此 \(\frac{\norm{Tx}_Y}{\norm{x}_X}=\norm{T(\frac{x}{\norm{x}_X})}_Y\). 当 \(\norm{x}_X=1\) 时, \(\frac{\norm{Tx}_Y}{\norm{x}_X}=\norm{Tx}_Y\).
        \end{rmk}

        \begin{ppt}
            {有界线性算子范数的性质}
            [Properties of Norm of Bounded Linear Operator]
            设 \(T\) 是有界线性算子, 则 \(\norm{Tx}\leq\norm{T}\norm{x}\).

        \end{ppt}

        \begin{ppt}
            {有界线性算子形成赋范线性空间}
            [Bounded Linear Operators Form a Normed Vector Space]
            设 \((X, \norm{\cdot}_X)\) 和 \((Y, \norm{\cdot}_Y)\) 是两个赋范线性空间. 定义 \(\blo(X,Y)\) 是从 \(X\) 到 \(Y\) 的所有有界线性算子组成的集合. 则 \(\blo(X,Y)\) 是赋范线性空间.
        \end{ppt}

        \begin{thm}
            {有界线性算子空间形成 Banach 空间的充要条件}
            设 \(X,Y\) 是两个赋范线性空间. 则 \(\blo(X,Y)\) 是 Banach 空间当且仅当 \(Y\) 是 Banach 空间.
        \end{thm}

        \begin{prf}
            设 \(\{T_n\}\) 是一列 \(X\) 到 \(Y\) 的有界线性算子, 使得 \(\{T_n\}\) 是 \(\blo(X,Y)\) 中的 Cauchy 列. 则对于任意 \(x\in X\), \(\{T_n x\}\) 是 \(Y\) 中的 Cauchy 列. 因为 \(Y\) 是完备的（Banach 空间）, 所以存在 \(y\in Y\) 使得 \(T_n x\to y\). 定义 \(T:X\to Y\) 使得 \(Tx=y\). \\\
            验证: (1) \(T\) 是有界线性算子: 线性性是显然的; 有界性是由于
            \[\norm{Tx}\leq\norm{T_N x}+\lim_{n\to\infty}\norm{T_n x-T_N x}\leq\norm{T_N x}+\epsilon\norm{x}\leq(\norm{T_N}+\epsilon)\norm{x}\]
            (2) \(T_n\) 依范数收敛到 \(T\): 
            \[\norm{T_n-T}=\sup_{\norm{x}=1}\norm{T_n x-Tx}=\sup_{\norm{x}=1}\lim_{k\to\infty}\norm{T_n x-T_k x}=0.\]
        \end{prf}

        \begin{crl}
            {依范数收敛蕴含范数收敛}
            若 \(T_n\to T\), 则 \(\norm{T_n}\to\norm{T}\) 作为实数序列收敛.
        \end{crl}

        \begin{crl}
            {依范数收敛蕴含逐点收敛}
            若 \(T_n\to T\), 则 \(\forall\,x\in X, T_n x\to Tx\).
        \end{crl}

        \begin{thm}
            {有界等价于连续}
            [Bounded iff Continuous]
            设 \(X,Y\) 是两个赋范线性空间, \(T:X\to Y\) 是一个线性映射. 则 \(T\) 是有界的当且仅当 \(T\) 是连续的.
        \end{thm}

        \begin{prf}
            在【数学分析III】中证明过该命题. 这里仅作提示.\\
            "(\(\implies\))": 设 \(T\) 是有界的, 则存在 \(M\geq 0\) 使得 \(\forall x\in X, \norm{Tx}_Y\leq M\norm{x}_X\). 则对于任意 \(\epsilon>0\), 当 \(\norm{x-x_0}_X<\frac{\epsilon}{M}\) 时, \(\norm{Tx-Tx_0}_Y=\norm{T(x-x_0)}_Y\leq M\norm{x-x_0}_X<\epsilon\). 从而 \(T\) 在 \(x_0\) 处连续. 因为 \(x_0\) 是任意的, 所以 \(T\) 是连续的.\\
            "(\(\impliedby\))": 设 \(T\) 是连续的, 则 \(T\) 在 \(0\) 处连续. 则存在 \(\delta>0\) 使得当 \(\norm{x}_X<\delta\) 时, \(\norm{Tx}_Y<1\). 根据赋范线性空间的平移不变性, 对于任意 \(x\in X\), 当 \(\norm{x}_X<\delta\) 时, \(\norm{Tx}_Y<1\). 从而对于任意 \(x\in X\), \(\norm{Tx}_Y<\frac{1}{\delta}\norm{x}_X\). 因此 \(T\) 是有界的.
        \end{prf}

        \begin{thm}
            {有界线性算子的复合是有界线性算子}
            [Composition of Bounded Linear Operators is Bounded]
            设 \(X,Y,Z\) 是三个赋范线性空间, \(T:X\to Y\) 和 \(S:Y\to Z\) 是有界线性算子. 则 \(S\circ T: X\to Z\) 也是有界线性算子. 且满足:
            \[\norm{S\circ T}\leq\norm{S}\norm{T}\]
        \end{thm}

        \begin{prf}
            令 \(\norm{x}=1\), 则
            \[\norm{S\circ T}=\sup\norm{S(T x)}=\sup\norm{S(y)}\leq\norm{S}\sup\norm{T x}=\norm{S}\norm{T}.\]
        \end{prf}

        \begin{thm}
            {有界线性算子的逆}
            设 $X$ 是 Banach 空间, $T \in \blo(X,X)$ 是双射, 则 $T^{-1} \in \blo (X,X)$.
        \end{thm}

        \begin{vardfn}
            {Banach 代数}
            [Banach Algebra]
            设 \(X\) 是一个 Banach 空间, 且 \(X\) 上定义了一个二元运算 \(\cdot\), 满足:
            \begin{enumerate}[label=(\roman*)]
                \item \((X, \cdot)\) 是代数, 满足代数运算公理;
                \item 有界性: \(\exists M\geq 0\) 使得 \(\forall x,y\in X, \norm{x\cdot y}\leq M\norm{x}\norm{y}\).
            \end{enumerate}
        \end{vardfn}

    \subsection{有限维 Banach 空间}

        \begin{intrormk}
            [Dedicatia]
            在有限维的情形下 Banach 空间的结构非常规整, 也有很多漂亮的性质.
        \end{intrormk}

        \begin{dfn}
            [DualSpace]
            {对偶空间}
            [Dual Space]
            设 \((X, \norm{\cdot})\) 是一个赋范线性空间. 定义 \(X\) 的\textbf{对偶空间} \(X^*\) 是 \(\blo(X, \R)\).
        \end{dfn}

        \begin{rmk}
            [Dedicatia]
            \(X^*\) 是 \(X\) 上的所有有界线性泛函组成的集合. 对偶空间中的元素都是线性泛函. \\
            需要注意的是, 这和作为线性空间的对偶空间的定义是不同的. 作为线性空间的对偶空间是所有线性泛函组成的集合, 而作为赋范线性空间的对偶空间是所有\textbf{有界线性泛函}组成的集合.\\
            只有在有限维线性空间中这两种定义才是相同的. 因为有限维的线性泛函一定是有界的.
        \end{rmk}

        \begin{thm}
            {有限维赋范线性空间的对偶空间同构于原空间}
            设 \(V\) 是赋范线性空间, \(\dim V<\infty\). 则 \(V^*\cong V\).
        \end{thm}

        \begin{prf}
            提示: 注意到二者的维数相同, 取定一组 \(V\) 的基 \(\{v_i\}_i^n\). 定义映射
            \[T:\qquad\begin{matrix}
                V & \to & V^* \\
                v_i & \mapsto & f_i
            \end{matrix}\]
            其中 \(f_i\) 是 \(V\) 上的线性泛函, 满足 \(f_i(v_j)=\delta_{ij}\). \\
            或者直接使用下面的同构定理.
        \end{prf}

        \begin{thm}
            {有限维赋范线性空间同构定理}
            设 \(V\) 是赋范线性空间, \(n=\dim V<\infty\). 则 \(V\cong\R^n\).
        \end{thm}

        \begin{rmk}
            [Dedicatia]
            这里我们说的同构指的是作为赋范线性空间的同构, 也就是说存在一个双射 \(T:V\to\R^n\) 使得 \(T\) 和 \(T^{-1}\) 都是有界线性算子. 
        \end{rmk}

        \begin{prf}
            取定一组 \(V\) 的基 \(\{v_i\}_i^n\). 定义映射
            \[T:\qquad\begin{matrix}
                \R^n & \to & V \\
                (x_1, x_2, \cdots, x_n) & \mapsto & \sum_{i=1}^n x_i v_i
            \end{matrix}\]
            验证:\\
            (1) 线性性显然成立;\\
            (2) \(T\) 是双射, 根据有限维线性空间的同构, 这也是显然的;\\
            (3) 考虑 \(T\) 的有界性:
            \[\norm{T(x_1, x_2, \cdots, x_n)}=\norm{\sum_{i=1}^n x_i v_i}\leq\sum_{i=1}^n|x_i|\norm{v_i}\leq\qty(\sum_{i=1}^n\abs{x_i})\qty(\sum_{i=1}^n\norm{v_i})  \]
            这说明
            \[\norm{T}\leq\sum_{i=1}^n\norm{v_i}<\infty\]
            (4) 考虑 \(T^{-1}\) 的有界性: 反证: 假如不是有界的, 那么不存在这样的 \(M\). 也就是说存在一列 \(\{x_n\}\) 使得
            \[\norm{T^{-1}(x_n)}>n\norm{x_n}\]
            根据齐次性, 则可以取定规范的 \(\{x_n\}\) s.t. \(\norm{x_n}=1\). 则 \(\norm{T^{-1}(x_n)}>n\). 令
            \[a_n = \frac{T^{-1}(x_n)}{\norm{T^{-1}(x_n)}}\]
            那么 \(\norm{a_n}=1\) 对所有 \(n\) 成立. 由于 \(V\) 是有限维的, 所以 \(V\) 中的单位球是紧的. 从而 \(\{a_n\}\) 中存在一个收敛子列 \(\{a_{n_k}\}\), 设其极限为 \(a\). 那么
            \[\norm{T^{-1}a} = \lim_{k \to \infty} \norm{T^{-1}a_{n_k}} = \lim_{k \to \infty} \frac{T^{-1}(x_{n_k})}{\norm{T^{-1}x_{n_k}}}\]
            由于 \(\norm{T^{-1}x_{n_k}}>n_k\), 所以 \(\norm{T^{-1}a}=0\). 从而 \(a=0\). 这与 \(\norm{a}=1\) 矛盾. 因此 \(T^{-1}\) 是有界的.
        \end{prf}

        \begin{thm}
            {有限维空间的范数等价}
            [Norms on Finite-Dimensional Spaces are Equivalent]
            设 \(V\) 是一个有限维的线性空间, \(\norm{\cdot}_1\) 和 \(\norm{\cdot}_2\) 是 \(V\) 上的两个范数. 则存在常数 \(C_1, C_2>0\) 使得
            \[\forall x\in V, C_1\norm{x}_1\leq\norm{x}_2\leq C_2\norm{x}_1\]
            或者等价地写作: 存在常数 \(C>0\) 使得
            \[\forall x\in V, \frac{1}{C}\norm{x}_1\leq\norm{x}_2\leq C\norm{x}_1.\]
        \end{thm}

        \begin{rmk}
            [Dedicatia]
            我们已经证明了决定有限维赋范线性空间的结构的唯一因素就是维数. 因此只要维度一样, 不论是什么范数一定可以建立同构. 自然也就是等价的.
        \end{rmk}

        \continued

\section{线性泛函分析基础}

    \subsection{Hahn-Banach 定理}

\end{document}